% 信号分析与处理 课程笔记 主文件
\documentclass[oneside,UTF8]{ctexbook}

% 核心宏包
\usepackage{amsmath,amssymb,amsfonts}
\usepackage{geometry}
\usepackage{graphicx}
\usepackage{booktabs}
\usepackage{longtable}
\usepackage{array}
\usepackage{enumitem}
\usepackage{tocbibind}
\usepackage{titlesec}

\usepackage{tcolorbox}
\tcbuselibrary{skins}

% 页面设置
\geometry{a4paper, left=2.5cm, right=2.5cm, top=2.5cm, bottom=2.5cm}

% 目录层级
\setcounter{tocdepth}{3}
\setcounter{secnumdepth}{3}

% 标题格式 —— 已修复反斜杠！
\titleformat{\chapter}[block]{\bfseries\Huge}{\thechapter}{1em}{}
\titleformat{\section}{\bfseries\Large}{\thesection}{1em}{}
\titleformat{\subsection}{\bfseries\large}{\thesubsection}{1em}{}
\titleformat{\subsubsection}{\bfseries}{\thesubsubsection}{1em}{}

% 抽样函数 Sa
\DeclareMathOperator{\Sa}{Sa}

% 超链接设置（蓝色无方框）
\usepackage{hyperref}
\hypersetup{
    colorlinks=true,
    linkcolor=blue,
    citecolor=blue,
    urlcolor=blue,
    pdfborder={0 0 0},
    pdfencoding=auto,
    pdfcreator=LaTeX
}

\begin{document}

% 封面文件 2.cover.tex
\thispagestyle{empty}
\restoregeometry
\newgeometry{left=2cm, right=2cm, top=2cm, bottom=2cm}

\begin{center}
    \vspace*{1cm}
    % 可替换为校徽/Logo
    % \includegraphics[width=0.6\textwidth]{hitsz_logo.png}
    % \vspace{1.5cm}
    
    {\Huge\bfseries 信号分析与处理}\\[1cm]
    {\Large\bfseries 课程笔记}\\[0.5cm]
    {\Large —— 时域/频域/复频域全解析 ——}\\[2cm]
    
    {\Large 整理者：Duckweed}\\[0.5cm]
    {\Large 日期：\today}\\[1cm]
    
    \vfill
    \rule{\linewidth}{0.5mm}
    \vspace{0.5cm}
    {\small 哈尔滨工业大学（深圳）}
\end{center}

\clearpage
\restoregeometry

\tableofcontents
\cleardoublepage
\setcounter{page}{1}

% ========== 绪论重点：信号的能量与功率分类（无错版） ==========
\chapter{绪论重点}

\section{信号的能量与功率分类（核心考点）}

\noindent 以下定义基于 \textbf{1$\Omega$ 电阻} 负载，瞬时功率 $P = |f(t)|^2$。

\subsection{1. 信号能量定义}
连续时间信号总能量：
\[
E = \lim_{T \to \infty} \int_{-T}^{T} |f(t)|^2 dt
\]
说明：$E$ 为信号总能量，单位 $\mathrm{J}$（焦耳）。

\subsection{2. 信号平均功率定义}
连续时间信号平均功率：
\[
P = \lim_{T \to \infty} \frac{1}{2T} \int_{-T}^{T} |f(t)|^2 dt
\]
说明：$P$ 为平均功率，单位 $\mathrm{W}$（瓦特）。

\subsection{3. 信号分类规则（正确教材定义）}
根据能量 $E$ 和功率 $P$ 的取值，信号分为三类：
\begin{itemize}
    \item \textbf{能量信号} \\
    判定：$E < \infty$（能量有限），此时必有 $P = 0$。\\
    示例：矩形脉冲、衰减指数信号。
    
    \item \textbf{功率信号} \\
    判定：$P < \infty$（功率有限），此时必有 $E = \infty$。\\
    示例：正弦信号、周期方波、单位阶跃信号。
    
    \item \textbf{非能量非功率信号} \\
    判定：\textbf{既不是能量信号，也不是功率信号} \\
    即：$E \to \infty$ 且 $P \to \infty$。\\
    示例：无界增长信号 $f(t)=t$、增长指数信号 $f(t)=e^{at}(a>0)$。
\end{itemize}

\begin{tcolorbox}[colback=gray!5!white,colframe=black!70!blue,title=考试必背结论]
\textbf{能量信号：} $E<\infty,\ \ P=0$ \\
\textbf{功率信号：} $P<\infty,\ \ E\to\infty$ \\
\textbf{非能量非功率：} 既不是能量信号，也不是功率信号
\end{tcolorbox}
% ========== 连续信号的时域描述和分析重点 ==========
% 已彻底修复 hyperref 警告：标题无数学公式、无特殊符号、纯文本书签
% 加粗正常渲染 + 优化排版布局 + 美观统一
\chapter{连续信号的时域描述和分析}

\section{正弦信号}
\begin{itemize}
\item 数学表达式：$f(t) = A\cos(\omega t + \varphi)$
\item 核心性质：
  \begin{itemize}
  \item \textbf{同频叠加性}：同频率正弦信号叠加后仍为同频率正弦信号。
  \item \textbf{异频合成}：周期比为有理数时为周期信号，否则非周期。
  \item \textbf{微积分不变性}：微分或积分后频率保持不变。
  \end{itemize}
\end{itemize}

\section{单位阶跃信号与矩形脉冲}
\begin{itemize}
\item 单位阶跃信号：
  \[
  u(t) = \begin{cases} 1, & t>0 \\ 0, & t<0 \end{cases}
  \]
\item 矩形脉冲信号：
  \[
  f_1(t)=u(t+1)-u(t-1)
  \]
  有效区间为 $\boldsymbol{[-1,1]}$，幅度为 $\boldsymbol{1}$，宽度为 $\boldsymbol{2}$。
\end{itemize}

\section{窗函数与信号截断}
\begin{itemize}
\item \textbf{矩形窗函数}：
  \[
  w(t) = u(t+a) - u(t-a)
  \]
  作用：保留区间 $\boldsymbol{[-a,a]}$ 内的信号，区间外信号强制置 $0$。

\item \textbf{信号相乘截断}：
  \[
  f_{\text{截断}}(t) = f(t) \cdot w(t)
  \]
  \begin{itemize}
  \item 窗内：信号保持不变
  \item 窗外：信号变为 $0$
  \end{itemize}

\item \textbf{典型例题结论}：
  \[
  \left[f_1(t)*f_2(t)\right]\cdot\left[u(t+5)-u(t-5)\right] = f_1(t+5)
  \]
  原因：
  \begin{enumerate}
  \item $f_1(t+5)$ 位于窗内，完整保留
  \item $f_1(t-5)$ 位于窗外，被完全截断
  \end{enumerate}
\end{itemize}

\section{单位冲激信号}
\begin{itemize}
\item 基本定义：
  \begin{itemize}
  \item 当 $t \neq 0$ 时，$\boldsymbol{\delta(t)=0}$
  \item 积分满足 $\boldsymbol{\displaystyle\int_{-\infty}^{\infty}\delta(t)dt=1}$
  \end{itemize}
\item \textbf{抽样特性（筛选特性）}：
  \[
  \int_{-\infty}^{\infty} f(\tau)\delta(\tau-t_0)d\tau = f(t_0)
  \]
\item \textbf{卷积性质}：
  \[
  f(t)*\delta(t-t_0)=f(t-t_0)
  \]
  信号与冲激卷积后，仅发生平移，形状保持不变。
\end{itemize}

\section{连续信号的时域运算}
\subsection{卷积运算}
\begin{itemize}
\item 定义式：
  \[
  f_1(t)*f_2(t)=\int_{-\infty}^{\infty}f_1(\tau)f_2(t-\tau)d\tau
  \]
\item 核心结论：
  \begin{itemize}
  \item \textbf{矩形脉冲与冲激卷积后仍为矩形脉冲。}
  \item \textbf{冲激只改变位置，不改变波形形状。}
  \end{itemize}
\end{itemize}

\subsection{信号时移与尺度变换}
\begin{itemize}
\item \textbf{平移变换}：$f(t\pm t_0)$
\item \textbf{尺度变换}：$f(at)$
\item \textbf{复合变换}：$y(0.5t-2)$ 包含平移与拉伸两步操作。
\end{itemize}

\section{正交函数集与完备性}
\begin{itemize}
\item \textbf{正交定义}：当 $k\neq l$ 时
  \[
  \int_{t_1}^{t_2}\phi_k(t)\phi_l(t)dt = 0
  \]

\item \textbf{完备性定义}：
若区间内 \textbf{任意信号} 都能被函数集线性组合无误差表示，则该集合为 \textbf{完备正交函数集}。

\item \textbf{正弦函数集性质}：
  \begin{itemize}
  \item 在区间 $\boldsymbol{(0,2\pi)}$ 上\textbf{正交}
  \item 在区间 $\boldsymbol{(0,\pi/2)}$ 上\textbf{不正交}
  \item 仅由正弦构成的集合\textbf{不完备}，只能表示\textbf{奇函数}，无法表示偶函数与直流信号
  \item 正弦 + 余弦集合才是\textbf{完备正交基}
  \end{itemize}
\end{itemize}

\section{冲激函数的抽样特性}
\subsection{典型积分公式}
\begin{enumerate}
    \item $\displaystyle \int_{-\infty}^{\infty} f(t-t_0)\delta(t)dt = \boldsymbol{f(-t_0)}$
    \item $\displaystyle \int_{-\infty}^{\infty} f(t_0-t)\delta(t)dt = \boldsymbol{f(t_0)}$
    \item $\displaystyle \int_{-\infty}^{\infty} \delta(t-t_0)u(t-2t_0)dt = \boldsymbol{u(-t_0)}$
    \item $\displaystyle \int_{-\infty}^{\infty} (t+\sin t)\delta\left(t-\frac{\pi}{6}\right)dt = \boldsymbol{\frac{\pi}{6}+\frac{1}{2}}$
    \item $\displaystyle \int_{-\infty}^{\infty} e^{-j\omega t}\left[\delta(t)-\delta(t-t_0)\right]dt = \boldsymbol{1 - e^{-j\omega t_0}}$
\end{enumerate}

\subsection{核心规则}
\begin{itemize}
    \item \textbf{冲激所在位置即为抽样点。}
    \item \textbf{将抽样点代入被积函数即可得到积分结果。}
    \item \textbf{被积函数在抽样点处必须连续。}
\end{itemize}

\section{周期信号判断方法}
\begin{itemize}
\item 各分量角频率：$\boldsymbol{\omega_1=10}$，$\boldsymbol{\omega_2=6}$
\item 基本周期：$\boldsymbol{T_1=\pi/5}$，$\boldsymbol{T_2=\pi/3}$
\item 周期比为有理数 $\boldsymbol{3/5}$，因此信号为周期信号。
\item 基波周期：$\boldsymbol{T=\pi}$。
\end{itemize}

\section{矩形脉冲与冲激的卷积}
\begin{itemize}
\item $\boldsymbol{f_1(t)}$ 为矩形脉冲。
\item $\boldsymbol{f_2(t)}$ 为两个冲激信号。
\item $\boldsymbol{s_1(t)}$ 与 $\boldsymbol{s_2(t)}$ 均为平移后的矩形脉冲组合。
\item \textbf{核心结论}：矩形脉冲与冲激卷积后仍然是矩形脉冲。
\end{itemize}

\section{信号卷积与波形变换}
\begin{itemize}
\item 分段卷积结果：
  \[
  y(t)=
  \begin{cases}
  0,&t\leq1\\
  2(t-1),&1<t<2\\
  6-2t,&2<t<3\\
  2t-8,&3<t<4\\
  0,&t\geq4
  \end{cases}
  \]
\item 波形变换 $\boldsymbol{y(0.5t-2)}$ 将区间 $\boldsymbol{(1,4)}$ 变换为 $\boldsymbol{(6,12)}$。
\end{itemize}

\section{核心结论总结}
\begin{tcolorbox}[colback=gray!5!white,colframe=black!70!blue,title=\textbf{考试必背结论}]
\begin{itemize}
    \item \textbf{冲激抽样}：代入冲激位置即可得到函数值。
    \item \textbf{冲激卷积}：信号仅平移，波形形状保持不变。
    \item \textbf{窗函数截断}：保留窗内信号，窗外置0。
    \item \textbf{完备性}：正弦+余弦完备，纯正弦不完备。
    \item \textbf{正弦集合}：在 $(0,2\pi)$ 上正交，但不完备。
    \item \textbf{矩形脉冲卷积冲激}：结果仍为矩形脉冲。
\end{itemize}
\end{tcolorbox}

\end{document}